% page 324 of the facsimile (image p-323 of the scan)
% block 154.9 x 254.9 mm ; left margin 8.0 mm
\pgc{-12.700mm}{0.000mm}{
\l{\cel{3}316}
\l{}
\l{\sou{kristalino}. (\sou{anat.)}\cel{3}Korpo lenso-forma, diafana, qua}
\l{\cel{3}esas en la okulo-globo, dop la pupilo, destinita a}
\l{\cel{3}refraktar la radii luma ed igar li konvergar a re-}
\l{\cel{3}tino. - DEFIRS.}
\l{}
\l{\sou{kristaloido}. (\sou{biol}.) Nomo di masi albuminoida ek subs-}
\l{\cel{3}tanco protoplasma, kun formi geometriala analoga ad}
\l{\cel{3}olti di la kristali.}
\l{}
\l{\sou{kriterio}. Karaktero per qua (per komparo) onu dicernas}
\l{\cel{3}lo exakta de lo ne-exakta, lo vera de lo falsa. DEFIRS.}
\l{}
\l{\sou{kriticismo}. (\sou{filo}z.)Doktrino qua igas diskutata la fidind-}
\l{\cel{3}eso di la idei di Raciono. - DEFIRS.}
\l{}
\l{\sou{kritikar}. (\sou{trans.}) I. Observar la qualesi, bona o mala, di}
\l{\cel{3}verko (literaturaja, arto-produkturo, e c.). - II. Igar}
\l{\cel{3}saliar la defekti di la kozi, di la personi. - DEFIRS.}
\l{}
\l{\sou{krizo}. I. Fazo danjeroza di maladeso-periodo, qua esos decid-}
\l{\cel{3}iganta, sive favoroze, sive perisige. II. Fazo danjeroza}
\l{\cel{3}quan transiras politiko, financo, industrio, komerco, e c.}
\l{\cel{3}- DEFIRS.}
\l{}
\l{\sou{krizalido}. Nimfo di lepidoptero e lua envelopilo. - DEFIS.}
\l{}
\l{\sou{krizantemo}.\cel{2}(\sou{bot.)}\cel{3}Arboreto qua formacas genero ek la}
\l{\cel{3}familio \textquotedbl{}kompozaji\textquotedbl{}, kultivata en la Europana gardeni}
\l{\cel{3}pro lua flori, flava, blanka, rozea,violea. - L.\cel{1}\sou{chrys}-}
\l{\cel{3}\sou{anthemum.}\cel{1}- DEFIRSL.}
\l{}
\l{krizokalko. - Aloyuro de zinko e kupro, kompozajo qua simu-}
\l{\cel{3}las oro e quan onu uzas en la faco di juveli falsa.-}
\l{\cel{3}DFIS.}
\l{}
\l{\sou{kroasar}. (\sou{netrans}.)(\sou{aludante la korvi, la rani)}.\cel{2}Agar la}
\l{\cel{3}krio qua karakterizas sua speco. F.}
\l{}
\l{\sou{krochar}. (\sou{trans.)}\cel{3}Facar objekti ek mashi, de lano, filo,}
\l{\cel{3}silko, e c., per mikra stango de fero o ligno, di qua}
\l{\cel{3}un pinto havas hoketo.EFR.}
\l{}
\l{\sou{kroketo}. Ludo-speco, de origino Angla, en qua onu pulsas}
\l{\cel{3}buli per frapi agita per ligna martelo, por igar ta buli}
\l{\cel{3}transsubirar arketi stekita en la sulo.- DEFIRS.}
\l{}
\l{\sou{krokodilo}. (\sou{zool.})\cel{2}Reptero sauria, amfibia, di la granda}
\l{\cel{3}fluvii di India e di Egiptia, timinda pro lua grandeso}
\l{\cel{3}e devoremeso. - L.\cel{1}\sou{crocodilus.}\cel{1}- DEFIRS.}
\l{}
\l{\sou{kromo}. (\sou{kemio)}.\cel{2}Korpo simpla, metala, di qua la kompozaji}
\l{\cel{3}esas remarkinda pro lia kolori. -\cel{1}\sou{Simb. kem.:}\cel{2}Kr.}
\l{}
\l{kromatiko. I. (\sou{muziko).}\cel{1}Qua konsistas ek intersequo de mi-}
\l{\cel{3}toni (mi-toni kromatika), kontraste kun la mi-toni dia-}
\l{\cel{3}tonika, plu mikra ye un komao. - II. Qua koncernas la}
\l{\cel{3}kolori. DEFIS.}
}
